%
% Capítulo 5
%
\chapter{Backend e API} \label{cap:backend}

O backend do sistema foi desenvolvido com o objetivo de centralizar a persistência remota de dados e disponibilizar uma API
capaz de suportar os principais fluxos funcionais da aplicação.
Esta componente desempenha um papel essencial na arquitetura da solução, funcionando como ponto de articulação entre a
aplicação Android e a base de dados PostgreSQL.

A introdução do backend permitiu ultrapassar as limitações de uma solução exclusivamente local, abrindo caminho para
uma estrutura mais próxima de um sistema real distribuído, onde os dados principais deixam de depender apenas do armazenamento
existente no dispositivo do utilizador.

\section{Escolha Tecnológica}

A implementação do backend foi realizada em Kotlin, utilizando Ktor como framework principal para criação da API REST.
Esta escolha revelou-se adequada ao contexto do projeto, sobretudo por permitir manter consistência com a linguagem
já utilizada no frontend e por oferecer uma base relativamente simples e leve para a construção do servidor.

A persistência foi suportada por PostgreSQL, escolhido por ser um sistema de gestão de bases de dados relacional robusto,
amplamente utilizado e adequado ao armazenamento estruturado das entidades do projeto.
Para a interação com a base de dados foi utilizada a biblioteca Exposed, que facilita a definição de tabelas e a execução
de operações de acesso a dados em Kotlin.
A gestão de conexões foi assegurada por HikariCP, permitindo configurar o acesso à base de dados de forma eficiente e consistente.

No conjunto, estas tecnologias permitiram construir uma base sólida para a versão beta, conciliando simplicidade de desenvolvimento
com capacidade de evolução futura.

\section{Estrutura do Backend}

O backend foi organizado em vários componentes com responsabilidades distintas, de forma a manter o código compreensível e modular.
Entre os principais elementos da sua estrutura destacam-se:
\begin{itemize}
    \item ficheiro principal de arranque da aplicação, responsável pela inicialização do servidor;
    \item configuração das rotas da API;
    \item modelos e DTOs utilizados na troca de informação entre cliente e servidor;
    \item serviços responsáveis por partes relevantes da lógica do sistema;
    \item configuração e inicialização da ligação à base de dados;
    \item definição das tabelas persistidas em PostgreSQL.
\end{itemize}

Esta organização permite separar as diferentes preocupações do sistema, distinguindo a exposição de endpoints, a representação dos dados,
a lógica funcional e o acesso à persistência.

\section{Modelo de Persistência no Backend}

A camada de persistência do backend foi concebida para suportar as entidades centrais do domínio do projeto.
Na versão beta, foram modeladas estruturas que permitem armazenar:
\begin{itemize}
    \item professores;
    \item alunos;
    \item associação entre alunos e professor;
    \item disponibilidades;
    \item restrições associadas ao professor.
\end{itemize}

A informação é persistida em PostgreSQL, permitindo manter um repositório central do estado do sistema.
Esta abordagem representa um passo importante relativamente às fases iniciais do projeto, em que a aplicação
dependia fortemente de persistência local. Com a introdução do backend, foi possível começar a transferir a
responsabilidade da gestão de dados principais para uma componente remota, mais adequada a uma solução multiutilizador.

No caso das disponibilidades, o backend armazena os intervalos registados por professores e alunos, utilizando essa informação
posteriormente para a criação de propostas de horário. No caso das restrições, os dados configurados pelo professor são também persistidos,
permitindo que a criação de horários seja baseada em regras efetivamente armazenadas e reutilizáveis.

\section{API REST Implementada}

A API desenvolvida no backend foi desenhada para suportar os fluxos principais do sistema. Os endpoints existentes permitem realizar operações
de registo, consulta, associação de utilizadores, gestão de disponibilidades, gestão de restrições e criação de horários.

Entre os principais endpoints implementados incluem-se:
\begin{itemize}
    \item verificação de disponibilidade do servidor;
    \item criação e consulta de professores;
    \item criação e consulta de alunos;
    \item associação e desassociação de alunos a professores;
    \item consulta de alunos de um professor;
    \item criação, remoção e consulta de disponibilidades;
    \item criação e consulta de restrições;
    \item criação de propostas de horário.
\end{itemize}

A utilização de uma API REST permitiu estabelecer um contrato claro entre o frontend e o backend, facilitando a integração progressiva da aplicação Android com a componente remota.
O formato JSON foi adotado como mecanismo de serialização dos dados, permitindo representar os pedidos e respostas de forma simples e consistente.

\section{Criação de Horários no Backend}

Uma das evoluções mais relevantes da versão beta foi a passagem da lógica de criação de horários para o backend.
Esta decisão permite centralizar a lógica principal do sistema e reduzir a duplicação de comportamento entre cliente e servidor.

A partir da informação armazenada, o backend:
\begin{itemize}
    \item obtém a disponibilidade do professor;
    \item obtém as disponibilidades dos alunos associados;
    \item carrega as restrições definidas;
    \item divide os intervalos em blocos temporais adequados;
    \item  cria sessões compatíveis com os dados disponíveis.
\end{itemize}

Esta lógica permite que a proposta de horário seja construída de forma coerente e reutilizável, garantindo que diferentes clientes possam consumir
o mesmo comportamento central sem necessidade de replicar o algoritmo localmente.

\section{Validação com Postman}

A validação do backend foi realizada com recurso ao Postman, ferramenta utilizada para testar os endpoints desenvolvidos antes da integração total com o frontend.
Este processo foi importante para confirmar que os contratos da API estavam corretos e que as operações principais produziam os resultados esperados.

Durante os testes foram validados fluxos como:
\begin{itemize}
    \item criação de professores e alunos;
    \item consulta de listas de utilizadores;
    \item associação entre aluno e professor;
    \item registo e consulta de disponibilidades;
    \item registo e consulta de restrições;
    \item criação de propostas de horário.
\end{itemize}

A utilização do Postman permitiu isolar o backend durante o processo de validação, facilitando a identificação de problemas antes da integração direta com a aplicação Android.
Este passo revelou-se particularmente útil para testar a coerência das respostas devolvidas pela API e a persistência correta dos dados em PostgreSQL.

\section{Integração com a Aplicação Android}

Após a validação inicial dos endpoints, a integração com o frontend foi sendo realizada de forma progressiva.
A aplicação Android passou a consumir a API para operações relacionadas com utilizadores, disponibilidades, restrições e, mais recentemente, criação de horários.

Esta integração permitiu aproximar a solução da sua arquitetura final, em que o backend assume o papel de componente central na gestão dos dados e da lógica principal do sistema.
Apesar de ainda existirem aspetos híbridos entre persistência local e remota, a evolução observada nesta fase demonstra que a base da comunicação cliente-servidor já se encontra estabelecida e funcional.

\section{Síntese do Estado Atual do Backend}

Na versão beta, o backend já se encontra suficientemente estruturado para suportar os principais fluxos do projeto.
A API implementada encontra-se funcional, a ligação à base de dados PostgreSQL foi estabelecida com sucesso e os principais
endpoints já foram validados de forma independente e em conjunto com a aplicação móvel.

Este estado representa um marco importante no desenvolvimento da solução, pois demonstra que o projeto já não depende apenas de lógica local no frontend.
O backend passou a constituir uma parte ativa e central do sistema, permitindo sustentar a evolução futura da solução para uma arquitetura mais consistente,
integrada e próxima da versão final pretendida.
